% GENERATED FILE. Edit canonical lessons, then run npm run generate:books.
% canonical-source-hash: fnv1a64:2439f769409a97fe
% canonical-lessons: RU-C38-rabotat, RU-C38-pet, RU-C38-tantsevat, RU-C38-smeyatsya, RU-C38-plakat

\chapter{Work, Sing, Dance, Laugh, Cry}
\label{ch:ru-work-sing-dance-laugh-cry}
\hlchaptermodality{38}

\begin{chapteropening}
\textbf{By the end of this chapter:} \emph{I can say to work, to sing, to dance, to laugh and to cry.}

Complete the last lesson of chapter 38: I can say to work, to sing, to dance, to laugh and to cry.
\end{chapteropening}

\section[\texorpdfstring{\ru{работать} (rabótat)}{работать (rabótat)}]{\ru{работать} (rabótat) — to work}

\label{lesson:RU-C38-rabotat}



\begin{quote}
Before the new one: say the Russian for a pupil, then the Russian for work.
\end{quote}

\subsection*{You'll want to know: \ru{работать}}
\textbf{\ru{работать}} (\emph{rabótat}) — \textquotedblleft{}to work\textquotedblright{}. The verb of \textbf{\ru{работа}}.

\begin{cousinweb}
\textbf{\ru{работать}} is \emph{\ru{работа}}, \textquotedblleft{}work\textquotedblright{}, made into a verb with \emph{-\ru{ать}}. \emph{\ru{работа}} goes back to a Proto-Slavic word for \textquotedblleft{}toil, slavery\textquotedblright{}, built on the root of \emph{\ru{раб}}, \textquotedblleft{}slave\textquotedblright{} (Vasmer).
\end{cousinweb}

\begin{grammarlens}[title={What you've built}]
One more word for this chapter.
\end{grammarlens}

\subsection*{Your turn}
\begin{itemize}
\raggedright
  \item \emph{Say it:} \emph{rabótat}
  \item \emph{Say it:} \emph{rabótat}, once more
  \item \emph{Say it:} say \emph{rabótat}
  \item \emph{Recall:} say the Russian for a pupil, then the Russian for work, then say \emph{rabótat} again
\end{itemize}

\begin{tcolorbox}[breakable,colback=teal!4,colframe=teal!35!black,title={Before you move on}]
What does \ru{работать} mean? (\textquotedblleft{}To work\textquotedblright{}.) Say it once more.
\end{tcolorbox}
\section[\texorpdfstring{\ru{петь} (pet)}{петь (pet)}]{\ru{петь} (pet) — to sing}

\label{lesson:RU-C38-pet}



\begin{quote}
Before the new one: say the Russian for work, then the Russian for to work.
\end{quote}

\subsection*{You'll want to know: \ru{петь}}
\textbf{\ru{петь}} (\emph{pet}) — \textquotedblleft{}to sing\textquotedblright{}.

\begin{cousinweb}
\textbf{\ru{петь}} $\leftarrow$ Proto-Slavic \emph{*pěti}, with \emph{*pojǫ}, \textquotedblleft{}I sing\textquotedblright{} — inherited. The deeper origin is unclear: Vasmer finds the proposed links unreliable, and Trubachev, in a note to Vasmer, explains it from \emph{\ru{пить}}, \textquotedblleft{}to drink\textquotedblright{}, since \emph{\ru{пою}} means both \textquotedblleft{}I sing\textquotedblright{} and \textquotedblleft{}I give to drink\textquotedblright{}.
\end{cousinweb}

\begin{grammarlens}[title={What you've built}]
One more word for this chapter.
\end{grammarlens}

\subsection*{Your turn}
\begin{itemize}
\raggedright
  \item \emph{Say it:} \emph{pet}
  \item \emph{Say it:} \emph{pet}, once more
  \item \emph{Say it:} say \emph{pet}
  \item \emph{Recall:} say the Russian for work, then the Russian for to work, then say \emph{pet} again
\end{itemize}

\begin{tcolorbox}[breakable,colback=teal!4,colframe=teal!35!black,title={Before you move on}]
What does \ru{петь} mean? (\textquotedblleft{}To sing\textquotedblright{}.) Say it once more.
\end{tcolorbox}
\section[\texorpdfstring{\ru{танцевать} (tantsevát)}{танцевать (tantsevát)}]{\ru{танцевать} (tantsevát) — to dance}

\label{lesson:RU-C38-tantsevat}



\begin{quote}
Before the new one: say the Russian for to work, then the Russian for to sing.
\end{quote}

\subsection*{You'll want to know: \ru{танцевать}}
\textbf{\ru{танцевать}} (\emph{tantsevát}) — \textquotedblleft{}to dance\textquotedblright{}.

\begin{cousinweb}
\textbf{\ru{танцевать}} is built on \emph{\ru{танец}}, \textquotedblleft{}a dance\textquotedblright{}, which came through Polish \emph{taniec} from Middle High German \emph{tanz}, \textquotedblleft{}a dance\textquotedblright{} (Vasmer).
\end{cousinweb}

\begin{grammarlens}[title={What you've built}]
One more word for this chapter.
\end{grammarlens}

\subsection*{Your turn}
\begin{itemize}
\raggedright
  \item \emph{Say it:} \emph{tantsevát}
  \item \emph{Say it:} \emph{tantsevát}, once more
  \item \emph{Say it:} say \emph{tantsevát}
  \item \emph{Recall:} say the Russian for to work, then the Russian for to sing, then say \emph{tantsevát} again
\end{itemize}

\begin{tcolorbox}[breakable,colback=teal!4,colframe=teal!35!black,title={Before you move on}]
What does \ru{танцевать} mean? (\textquotedblleft{}To dance\textquotedblright{}.) Say it once more.
\end{tcolorbox}
\section[\texorpdfstring{\ru{смеяться} (smeyátsya)}{смеяться (smeyátsya)}]{\ru{смеяться} (smeyátsya) — to laugh}

\label{lesson:RU-C38-smeyatsya}



\begin{quote}
Before the new one: say the Russian for to sing, then the Russian for to dance.
\end{quote}

\subsection*{You'll want to know: \ru{смеяться}}
\textbf{\ru{смеяться}} (\emph{smeyátsya}) — \textquotedblleft{}to laugh\textquotedblright{}.

\begin{cousinweb}
\textbf{\ru{смеяться}} $\leftarrow$ Proto-Slavic \emph{*smĭjati sę} — inherited; \emph{-\ru{ся}} is the old \textquotedblleft{}oneself\textquotedblright{}. Vasmer relates it to Latvian \emph{smiêt}, \textquotedblleft{}to laugh\textquotedblright{}, Sanskrit \emph{smáyate}, \textquotedblleft{}smiles\textquotedblright{}, Greek \emph{meidiáō}, \textquotedblleft{}I smile\textquotedblright{}, and Middle English \emph{smilen}, so English \emph{smile} is a cousin.
\end{cousinweb}

\begin{grammarlens}[title={What you've built}]
One more word for this chapter.
\end{grammarlens}

\subsection*{Your turn}
\begin{itemize}
\raggedright
  \item \emph{Say it:} \emph{smeyátsya}
  \item \emph{Say it:} \emph{smeyátsya}, once more
  \item \emph{Say it:} say \emph{smeyátsya}
  \item \emph{Recall:} say the Russian for to sing, then the Russian for to dance, then say \emph{smeyátsya} again
\end{itemize}

\begin{tcolorbox}[breakable,colback=teal!4,colframe=teal!35!black,title={Before you move on}]
What does \ru{смеяться} mean? (\textquotedblleft{}To laugh\textquotedblright{}.) Say it once more.
\end{tcolorbox}
\section[\texorpdfstring{\ru{плакать} (plákat)}{плакать (plákat)}]{\ru{плакать} (plákat) — to cry}

\label{lesson:RU-C38-plakat}



\begin{quote}
Before the new one: say the Russian for to dance, then the Russian for to laugh.
\end{quote}

\subsection*{You'll want to know: \ru{плакать}}
\textbf{\ru{плакать}} (\emph{plákat}) — \textquotedblleft{}to cry\textquotedblright{}.

\begin{cousinweb}
\textbf{\ru{плакать}} $\leftarrow$ Proto-Slavic \emph{*plakati} — inherited. Vasmer gives the first Slavic sense as \textquotedblleft{}to beat one's breast\textquotedblright{} in mourning, and relates it to Lithuanian \emph{plàkti}, \textquotedblleft{}to beat\textquotedblright{}, Greek \emph{plēssō}, \textquotedblleft{}I strike\textquotedblright{}, and Latin \emph{plangō}, \textquotedblleft{}I beat my breast, lament\textquotedblright{}.
\end{cousinweb}

\begin{grammarlens}[title={What you've built}]
One more word for this chapter.
\end{grammarlens}

\subsection*{Your turn}
\begin{itemize}
\raggedright
  \item \emph{Say it:} \emph{plákat}
  \item \emph{Say it:} \emph{plákat}, once more
  \item \emph{Say it:} say \emph{plákat}
  \item \emph{Recall:} say the Russian for to dance, then the Russian for to laugh, then say \emph{plákat} again
\end{itemize}

\begin{tcolorbox}[breakable,colback=teal!4,colframe=teal!35!black,title={Before you move on}]
What does \ru{плакать} mean? (\textquotedblleft{}To cry\textquotedblright{}.) Say it once more.
\end{tcolorbox}
